% 03_HumanGate卡.tex —— 作业要求③ ｜ Human Gate 卡 × 5 ｜ A3 横版
% 编译：xelatex -jobname=gatecards 03_HumanGate卡.tex
\documentclass[10pt]{article}
\usepackage[a3paper,landscape,margin=10mm]{geometry}
\usepackage{xeCJK}
\setCJKmainfont{PingFang SC}
\setCJKsansfont{PingFang SC}
\xeCJKDeclareCharClass{CJK}{"2460 -> "2473, "2265, "2264, "2192}
\usepackage{array}
\usepackage[table]{xcolor}
\usepackage{colortbl}
\pagestyle{empty}
\setlength{\parindent}{0pt}
\setlength{\parskip}{0pt}

\definecolor{ink}{HTML}{1B2A38}
\definecolor{headbg}{HTML}{19324D}
\definecolor{hair}{HTML}{C9D6E2}
\definecolor{gedge}{HTML}{C2701A}
\definecolor{gtint}{HTML}{FDF1E0}
\definecolor{aiedge}{HTML}{0F6E80}
\definecolor{aitint}{HTML}{EDF5F8}
\definecolor{ph}{HTML}{6B4FA8}          % 占位符颜色
\definecolor{muted}{HTML}{5A6B7B}

\newcommand{\cardtitle}[1]{{\fontsize{10}{12}\selectfont\bfseries\textcolor{white}{#1}}}
\newcommand{\flab}[1]{{\fontsize{8}{9.8}\selectfont\bfseries\textcolor{gedge}{#1}}}
\newcommand{\fbody}[1]{{\fontsize{7.7}{9.4}\selectfont #1}}
\newcommand{\phv}[1]{{\fontsize{7.7}{9.4}\selectfont\bfseries\textcolor{ph}{#1}}}
\newcommand{\back}[1]{{\fontsize{7.9}{9.6}\selectfont\bfseries\textcolor{ink}{#1}}}

% 一张卡 = 单列表格，5 行
\newcommand{\gatecard}[6]{%
\begin{minipage}[t]{12.4cm}
\renewcommand{\arraystretch}{1.42}
\setlength{\tabcolsep}{5pt}
\begin{tabular}{@{}>{\raggedright\arraybackslash}p{12.4cm}@{}}
\rowcolor{gedge}\cardtitle{#1}\\[0.4mm]
\rowcolor{gtint}\flab{AI 提供什么}\quad\fbody{#2}\\[0.4mm]
\rowcolor{white}\flab{看什么}\quad\fbody{#3}\\[0.4mm]
\rowcolor{aitint}\flab{能做什么}\quad\fbody{#4}\\[0.4mm]
\rowcolor{gtint}\back{不通过回哪}\quad\back{#5}\\
\end{tabular}
\vspace{1.6mm}

{\fontsize{6.9}{8.4}\selectfont\textcolor{muted}{旁注（填写示例）：#6}}
\end{minipage}}

\begin{document}

{\fontsize{22}{26}\selectfont\bfseries\textcolor{ink}{社区服务信息主视觉设计 · Human Gate 卡（5 张）}}\\[1.5mm]
{\fontsize{10}{12}\selectfont\textcolor{muted}{每卡 5 个字段：Gate 名称 / AI 提供什么 / 看什么 / 能做什么 / 不通过回哪 ｜ 张数 = 工作流程图里 Human Gate 的数量（G1–G5）｜ \textcolor{ph}{紫色粗体}=待人类提供的占位符；风格先按\textbf{五大类}给（references/05-style-spec.md），15 类暂不细分}}\\[4mm]

\gatecard{G1 · 分类与路由确认}%
{theme.yaml（五轴取值 + 唯一主骨架 + language）；分类依据说明；缺类与越界项清单}%
{领域是否分对（D01–D15）／骨架是否匹配信息功能（S1–S9）／受众有无遗漏（老年人·残障·国际居民）／权威等级有无高估／语言版本是否与受众一致}%
{通过 1 个 ／ 改判领域或骨架 ／ 调整受众或权威等级 ／ 全部否决退回重分类}%
{回「主题分类与路由」（步骤 2 / N2）}%
{「受众漏了国际居民」→ 判定语言版本应为双语，退回 N2 重判}%
\hfill
\gatecard{G2 · 事实与母本校对}%
{fact-table.md（逐条字段 + 官方来源）；copy-diff.md（母本原文 vs 成品实际文本）；source.md（链接 / 文号）}%
{字符级一致性（标点·空格·全半角·连字符）／电话与时间是否逐位正确／来源是否可追溯／有无 AI 润色或新增承诺}%
{通过 ／ 标出并退回改正 ／ 要求补来源 ／ 判定母本本身缺失（升级为补料）}%
{回「事实结构化」（步骤 3 / N3）；\textbf{母本本身缺失 → 回「接收输入包」（步骤 1 / N1）}}%
{「电话少一位」→ 标红退回 N3，并在 needs-human 登记待补}%
\hfill
\gatecard{G3 · 设计系统冻结}%
{candidates.md（≥3 个可区分方向，各含关键词·构图逻辑·风格风险）；每方向的色卡 / 字体 / 栅格建议；\phv{本主题所属大类的 \{\{F*\_KEYWORDS\}\}、\{\{F*\_COLORS\}\}、\{\{F*\_FONTS\}\}、\{\{F*\_MUST\}\}、\{\{F*\_BAN\}\}（见 references/05-style-spec.md）}}%
{目标匹配 / 信息清晰 / 品牌适配 / 视觉潜力（课件四维度）＋ 是否符合 \phv{\{\{F*\_KEYWORDS\}\}} ＋ 色与字体是否落在 \phv{\{\{F*\_COLORS\}\}} / \phv{\{\{F*\_FONTS\}\}} 之内 ＋ \phv{\{\{F*\_MUST\}\}} 是否齐备 ＋ 是否触碰 \phv{\{\{F*\_BAN\}\}} 的禁令}%
{通过 1 个 ／ 合并 2 个 ／ 全部否决退回补候选 ／ 指定以某张参考图为基准重出}%
{回「方向候选生成」（步骤 4 / N4）}%
{「方向 2 的配色接近参考图，但字体太活泼」→ 合并方向 2+3 并要求换字体}%
\vspace{4mm}

\gatecard{G4 · 内容与可读性确认}%
{poster 成稿（A3 画布，PNG + PDF）；check-report.md（字号·对比度实测值）；\phv{\{\{F*\_MUST\}\} 落实对照 + \{\{STYLE\_EXEC\_REFS\}\}}}%
{信息项是否 100\% 完整／电话与时间不跨行、无遮挡／字号 ≥26 pt（含老年人·残障 ≥31 pt）／对比度 ≥4.5:1（升档 ≥7:1）／安全边距 ≥18 mm／有无禁止元素}%
{通过 ／ 退回改排版 ／ 退回改骨架 ／ 要求补素材}%
{回「中文版入版排版」（步骤 7 / N7）；\textbf{骨架问题 → 回「主视觉线框」（步骤 5 / N5）}}%
{「对比度实测 4.1:1」→ 附实测截图，退回 N7 调色}%
\hfill
\gatecard{G5 · 双语一致性与对外责任}%
{中英成稿（或 zh-only 说明）；派生物料确认记录；落款·发布主体·时效字段清单}%
{双语是否等价、有无删词／语言版本是否符合受众规则（含 A06 才出双语）／落款与发布主体是否可对外／发布时间与有效期是否齐全（T4 必填）／免责口径是否足够}%
{通过 ／ 退回改英文重排 ／ 退回改派生物料 ／ 要求补落款与时效}%
{回「英文版本地化重排」（步骤 8 / N8）；\textbf{派生物料问题 → 回「派生物料询问与生成」（步骤 9 / N9）}}%
{「英文删了“定期”」→ 违反正文零删词，退回 N8 重排}%
\hfill
\begin{minipage}[t]{12.4cm}
\renewcommand{\arraystretch}{1.3}
\setlength{\tabcolsep}{4pt}
\begin{tabular}{@{}>{\raggedright\arraybackslash}p{12.4cm}@{}}
\rowcolor{headbg}{\fontsize{9}{11}\selectfont\bfseries\textcolor{white}{填写规则（5 字段都不许写“人工确认”四个字）}}\\[0.4mm]
\rowcolor{white}\fbody{\textbf{① Gate 名称}：它决定什么，写成“决定……”的短语，不写“审核”两字。}\\[0.3mm]
\rowcolor{aitint}\fbody{\textbf{② AI 提供什么}：人拿到手就能判断的材料，必须是**具体文件或具体清单**。}\\[0.3mm]
\rowcolor{white}\fbody{\textbf{③ 看什么}：判断维度，写 3–5 个，用“／”分隔；能对应脚本检查的要写阈值。}\\[0.3mm]
\rowcolor{aitint}\fbody{\textbf{④ 能做什么}：**枚举**允许动作（3–4 个），不允许出现“视情况而定”。}\\[0.3mm]
\rowcolor{white}\fbody{\textbf{⑤ 不通过回哪}：必须写**具体节点名 + 步骤号**，只写“退回重做”不合格。}\\[0.3mm]
\rowcolor{gtint}\fbody{\textbf{占位符}：\phv{\{\{...\}\}} 一律紫色粗体；人类给资料后按下方清单逐项替换，替换即视为一次冻结变更。}\\
\end{tabular}
\end{minipage}

\vspace{5mm}
{\fontsize{11}{13}\selectfont\bfseries\textcolor{ink}{待人类提供的风格与参考资料清单（占位符索引）}}\\[2mm]
\renewcommand{\arraystretch}{1.5}
\setlength{\tabcolsep}{4pt}
\begin{tabular}{@{}>{\raggedright\arraybackslash}p{4.4cm}%
>{\raggedright\arraybackslash}p{7.6cm}%
>{\raggedright\arraybackslash}p{6.6cm}%
>{\centering\arraybackslash}p{2.4cm}%
>{\centering\arraybackslash}p{2.4cm}%
>{\raggedright\arraybackslash}p{14.6cm}@{}}
\rowcolor{headbg}
{\fontsize{8.4}{10}\selectfont\bfseries\textcolor{white}{占位符}} &
{\fontsize{8.4}{10}\selectfont\bfseries\textcolor{white}{要什么}} &
{\fontsize{8.4}{10}\selectfont\bfseries\textcolor{white}{格式要求}} &
{\fontsize{8.4}{10}\selectfont\bfseries\textcolor{white}{用在哪张卡}} &
{\fontsize{8.4}{10}\selectfont\bfseries\textcolor{white}{什么时候要}} &
{\fontsize{8.4}{10}\selectfont\bfseries\textcolor{white}{给了之后怎么处理}} \\
\hline
\phv{\{\{F1–F5\_KEYWORDS\}\}} & \fbody{3 个风格关键词，每词附“意味着什么 / 不意味着什么”（每大类一份）} & \fbody{纯文本，每词 1–2 句} & \fbody{G3} & \fbody{冻结前} & \fbody{写入设计系统的「风格关键词」节，转成可校验规则（AI 出候选时逐条对照）} \\
\phv{\{\{MOODBOARD\}\}} & \fbody{3–5 张参考图 / 情绪板} & \fbody{PNG / JPG，长边 ≥1500 px；注明来源与可否商用} & \fbody{G3} & \fbody{冻结前} & \fbody{逐张拆解：构图 / 色彩 / 字体 / 留白 / 图形语言；输出“可复用规则 + 不可照搬元素”两张清单} \\
\phv{\{\{F1–F5\_COLORS\}\}} & \fbody{主色·辅助色·警示色（\#RRGGBB + 允许/禁止用途）} & \fbody{Hex 列表或色卡图} & \fbody{G3} & \fbody{冻结前} & \fbody{生成色卡表（色值 + 用途 + 对比度实测值），并标出允许 / 禁止用法} \\
\phv{\{\{F1–F5\_FONTS\}\}} & \fbody{中英文字体（名字 + 字重 + 授权）} & \fbody{字体名 + 样张 + 授权说明} & \fbody{G3} & \fbody{冻结前} & \fbody{建字号阶梯表（pt/mm）＋ 回退字族 ＋ 商用授权核对} \\
\phv{\{\{F1–F5\_MUST\}\}} & \fbody{必须包含的元素（可逐条校验的清单）} & \fbody{清单，每项可判定} & \fbody{G3} & \fbody{冻结前} & \fbody{转成校验脚本检查项（第 17–18 项），并在 Gate 4 逐条打勾} \\
\phv{\{\{F1–F5\_BAN\}\}} & \fbody{明确禁止的元素（可判定的禁令）} & \fbody{清单，看到即可判违规} & \fbody{G3} & \fbody{冻结前} & \fbody{写入反模式清单与禁止元素扫描；触碰即 Gate 3 否决} \\
\phv{\{\{STYLE\_EXEC\_REFS\}\}} & \fbody{成品与参考图的对照样本} & \fbody{图：成品局部 + 对应参考图} & \fbody{G4} & \fbody{出稿后} & \fbody{做落地偏差分析：哪些还原、哪些走样、逐条给修改建议} \\
\hline
\end{tabular}

\vspace{4mm}
{\fontsize{8}{10}\selectfont\textcolor{muted}{说明：① 卡数 = 工作流程图 Human Gate 数（G1–G5，共 5 张，满足课件“至少 3 张”）；② 每卡的“不通过回哪”与流程图橙色虚线一一对应：G1→2、G2→3、G3→4、G4→7（骨架回 5）、G5→8（派生回 9）；③ 风格占位符以 G3 为主、G4 为辅（G4 只收“落地效果对照”一类）；④ 人类提供资料 → 替换占位符 → 视为一次冻结变更，须记入 CHANGELOG 并重跑一次迷你 Gate 3。}}

\end{document}
